\documentclass[10pt]{extarticle}
\usepackage[a4paper,margin=1.5cm]{geometry}
\usepackage[T1]{fontenc}
\usepackage{amsmath,amssymb}
\usepackage{graphicx}
\usepackage[export]{adjustbox}
\usepackage{multicol}
\usepackage{enumitem}
\usepackage{booktabs,array,tabularx}
\setlength{\columnsep}{0.7cm}
\setlength{\parindent}{0pt}
\setlength{\parskip}{3pt}
\pagestyle{empty}
\makeatletter
% figure registry (figdefs.tex is generated by the extract step)
\newcommand{\qfigdef}[4]{\@namedef{qf@file@#1}{#2}\@namedef{qf@w@#1}{#3}\@namedef{qf@wide@#1}{#4}}
\newcommand{\qf@check}[1]{\@ifundefined{qf@file@#1}{\PackageError{score2}{Unknown figure id #1}{}}{}}
\newcommand{\qf@img}[4][]{\includegraphics[width=#2\dimexpr\@nameuse{qf@w@#4}pt\relax,max width=#3,max height=0.45\textheight,#1]{\@nameuse{qf@file@#4}}}
% Full-width questions: finalize wraps every question that contains a wide figure in \qfullbegin...\qfullend,
% so the WHOLE question (stem, figures, options) is set across the page in reading order, then the two columns
% resume. The question's head - its number, stem and first figure (\qfullsplit marks the end of the head) -
% is kept on one page, so a number is never left at a page foot with its figure overleaf; the rest may break
% across pages normally, so no big gaps are left. Without a \qfullsplit, the whole question is kept together
% if it is under 45% of a page. Never write these yourself.
\newif\ifqf@full
\newif\ifqf@boxed
\newsavebox{\qf@box}
\newcommand{\qf@open}{\global\qf@fulltrue\par\global\qf@boxedtrue
  \setbox\qf@box\vbox\bgroup\hsize\textwidth\linewidth\textwidth}
\newcommand{\qf@emit}[1]{\par\egroup\global\qf@boxedfalse
  \ifdim\dimexpr\ht\qf@box+\dp\qf@box\relax<#1\textheight\noindent\usebox{\qf@box}\par\else\unvbox\qf@box\fi}
\newcommand{\qfullsplit}{\ifqf@boxed\qf@emit{0.85}\fi}
\newcommand{\qf@close}{\ifqf@boxed\qf@emit{0.45}\fi\global\qf@fullfalse}
\newcommand{\qfullbegin}{\end{multicols}\qf@open}
\newcommand{\qfullend}{\qf@close\begin{multicols}{2}}
\newcommand{\qfullstart}{\qf@open}
\newcommand{\qfulldone}{\qf@close}
% \qfig[scale]{id}: block figure, centred, directly where it is written (size set by the tool).
% Fallback if a wide figure is not inside a full-width question: break out of the columns in place - no floats.
\newcommand{\qfig}[2][1]{\qf@check{#2}%
  \ifqf@full\qf@blk{#1}{#2}\else
  \ifnum\@nameuse{qf@wide@#2}=1 \expandafter\@firstoftwo\else\expandafter\@secondoftwo\fi
  {\end{multicols}\begin{center}\qf@img{#1}{\textwidth}{#2}\end{center}\begin{multicols}{2}}%
  {\qf@blk{#1}{#2}}\fi}
\newcommand{\qf@blk}[2]{\par\smallskip{\centering\qf@img{#1}{\linewidth}{#2}\par}\smallskip}
% \qfigi[scale]{id}: inline figure (inside option lists / tables / minipages); never breaks columns.
\newcommand{\qfigi}[2][1]{\qf@check{#2}\qf@img[valign=c]{#1}{\linewidth}{#2}}
\makeatother
% structure macros - lint reads these to check completeness and figure placement
\newcommand{\papersec}[1]{\par\bigskip{\centering\Large\bfseries PAPER-#1\par}\medskip}
\newcommand{\partsec}[2]{\par\medskip\noindent{\large\bfseries PART-#1: #2}\par\smallskip}
\newcommand{\sect}[1]{\par\smallskip\noindent\textbf{#1}\par}
\newcommand{\qn}[1]{\par\smallskip\noindent\textbf{#1.}~}
\newcommand{\paraq}[2]{\par\smallskip\noindent\textbf{Paragraph for Questions #1 and #2}\par}
\newenvironment{qoptions}{\begin{itemize}[leftmargin=1.8em,itemsep=1pt,topsep=2pt,label={}]}{\end{itemize}}
\newcommand{\opt}[2]{\item[(#1)] #2}
% figure registry (generated by score2_tools.py finalize)
\qfigdef{p1_phy_s0_q01_1}{SRGBT - Test 11 (03 Oct 2021)-img/p1_phy_s0_q01_1.png}{103.2}{0}
\qfigdef{p1_phy_s0_q02_1}{SRGBT - Test 11 (03 Oct 2021)-img/p1_phy_s0_q02_1.png}{61.3}{0}
\qfigdef{p1_phy_s0_q03_1}{SRGBT - Test 11 (03 Oct 2021)-img/p1_phy_s0_q03_1.png}{94.1}{0}
\qfigdef{p1_phy_s0_q04_2}{SRGBT - Test 11 (03 Oct 2021)-img/p1_phy_s0_q04_2.png}{78.9}{0}
\qfigdef{p1_phy_s0_q05_1}{SRGBT - Test 11 (03 Oct 2021)-img/p1_phy_s0_q05_1.png}{69.3}{0}
\qfigdef{p1_phy_s0_q06_1}{SRGBT - Test 11 (03 Oct 2021)-img/p1_phy_s0_q06_1.png}{96.7}{0}
\qfigdef{p1_phy_s0_q07_1}{SRGBT - Test 11 (03 Oct 2021)-img/p1_phy_s0_q07_1.png}{125.7}{0}
\qfigdef{p1_phy_s0_q09_1}{SRGBT - Test 11 (03 Oct 2021)-img/p1_phy_s0_q09_1.png}{80.4}{0}
\qfigdef{p1_phy_s0_q10_1}{SRGBT - Test 11 (03 Oct 2021)-img/p1_phy_s0_q10_1.png}{81.2}{0}
\qfigdef{p1_phy_s0_q11_1}{SRGBT - Test 11 (03 Oct 2021)-img/p1_phy_s0_q11_1.png}{60.9}{0}
\qfigdef{p1_phy_s0_q12_2}{SRGBT - Test 11 (03 Oct 2021)-img/p1_phy_s0_q12_2.png}{207.7}{0}
\qfigdef{p1_phy_s0_q13_1}{SRGBT - Test 11 (03 Oct 2021)-img/p1_phy_s0_q13_1.png}{196.1}{0}
\qfigdef{p1_phy_s0_q14_1}{SRGBT - Test 11 (03 Oct 2021)-img/p1_phy_s0_q14_1.png}{43.0}{0}
\qfigdef{p1_phy_s0_q15_1}{SRGBT - Test 11 (03 Oct 2021)-img/p1_phy_s0_q15_1.png}{164.2}{0}
\qfigdef{p1_phy_s0_q16_1}{SRGBT - Test 11 (03 Oct 2021)-img/p1_phy_s0_q16_1.png}{119.5}{0}
\qfigdef{p1_phy_s0_q18_1}{SRGBT - Test 11 (03 Oct 2021)-img/p1_phy_s0_q18_1.png}{104.3}{0}
\qfigdef{p1_chem_s0_q19_1}{SRGBT - Test 11 (03 Oct 2021)-img/p1_chem_s0_q19_1.png}{120.5}{0}
\qfigdef{p1_chem_s0_q19_2}{SRGBT - Test 11 (03 Oct 2021)-img/p1_chem_s0_q19_2.png}{70.4}{0}
\qfigdef{p1_chem_s0_q19_3}{SRGBT - Test 11 (03 Oct 2021)-img/p1_chem_s0_q19_3.png}{71.8}{0}
\qfigdef{p1_chem_s0_q19_4}{SRGBT - Test 11 (03 Oct 2021)-img/p1_chem_s0_q19_4.png}{70.4}{0}
\qfigdef{p1_chem_s0_q19_5}{SRGBT - Test 11 (03 Oct 2021)-img/p1_chem_s0_q19_5.png}{67.7}{0}
\qfigdef{p1_chem_s0_q30_1}{SRGBT - Test 11 (03 Oct 2021)-img/p1_chem_s0_q30_1.png}{76.5}{0}
\qfigdef{p1_chem_s0_q32_1}{SRGBT - Test 11 (03 Oct 2021)-img/p1_chem_s0_q32_1.png}{95.8}{0}
\qfigdef{p1_chem_s0_q33_1}{SRGBT - Test 11 (03 Oct 2021)-img/p1_chem_s0_q33_1.png}{180.6}{0}
\qfigdef{p1_chem_s0_q36_1}{SRGBT - Test 11 (03 Oct 2021)-img/p1_chem_s0_q36_1.png}{149.4}{0}
\raggedcolumns
\makeatletter\@beginparpenalty=10000\makeatother
\newcolumntype{L}[1]{>{\raggedright\arraybackslash}p{#1}}
\begin{document}
\begin{center}{\Large\bfseries SRGBT --- Test 11 (03 Oct 2021)}\end{center}
\begin{multicols}{2}
\raggedcolumns
\papersec{1}
\partsec{1}{PHYSICS}
\qn{1} Two blocks of mass m$_1$ and m$_2$ are connected by a spring of force constant k. Block of mass m$_1$ is pulled by a constant force F$_1$ and other block is pulled by a constant force F$_2$. Find the maximum elongation in the spring.\par
\qfig{p1_phy_s0_q01_1}
\begin{qoptions}
\opt{A}{$\left[\dfrac{\mathrm F_1\mathrm m_2+\mathrm m_1\mathrm F_2}{\mathrm m_1+\mathrm m_2}\right]\dfrac1{\mathrm k}$}
\opt{B}{$\left[\dfrac{\mathrm F_1\mathrm m_2+\mathrm m_1\mathrm F_2}{\mathrm m_1+\mathrm m_2}\right]\dfrac2{\mathrm k}$}
\opt{C}{$\left[\dfrac{\mathrm F_1\mathrm m_2+\mathrm m_1\mathrm F_2}{\mathrm m_1+\mathrm m_2}\right]\dfrac3{\mathrm k}$}
\opt{D}{$\left[\dfrac{\mathrm F_1\mathrm m_2+\mathrm m_1\mathrm F_2}{\mathrm m_1+\mathrm m_2}\right]\dfrac4{\mathrm k}$}
\end{qoptions}
\qn{2} A chain of length L and mass per unit length $\rho$ is pilled on a horizontal surface. One end of the chain is lifted vertically with a constant velocity $v$ by a variable force P. Which of the following is \textbf{incorrect}.\par
\qfig{p1_phy_s0_q02_1}
\begin{qoptions}
\opt{A}{Force P at height L/2 is $\rho\left[g\dfrac L2+\mathrm V^2\right]$}
\opt{B}{Total work done by force P during lifting the chain is $\dfrac{\rho gL^2}2+\rho\mathrm V^2\mathrm L$}
\opt{C}{Energy lost during lifting the chain is $\rho\mathrm{LV}^2$}
\opt{D}{Energy lost during lifting the chain is $\dfrac12\rho\mathrm{LV}^2$}
\end{qoptions}
\qn{3} Five identical balls each of mass m and radius r are strung like beads at random and are at rest along a smooth, rigid horizontal thin rod of length L, mounted between immovable supports as shown in the figure. Assume 10r $<$ L and that the collision between balls or between balls and supports are elastic. If one ball is struck horizontally so as to acquire a speed v, the average force felt by the support is\par
\qfig{p1_phy_s0_q03_1}
\begin{qoptions}
\opt{A}{$\dfrac{5\mathrm{mv}^2}{\mathrm L-5\mathrm r}$}
\opt{B}{$\dfrac{\mathrm{mv}^2}{\mathrm L-10\mathrm r}$}
\opt{C}{$\dfrac{5\mathrm{mv}^2}{\mathrm L-10\mathrm r}$}
\opt{D}{$\dfrac{\mathrm{mv}^2}{\mathrm L-5\mathrm r}$}
\end{qoptions}
\qn{4} A is a fixed point at a height h above a perfectly inelastic smooth surface. A light inextensible string of length $\ell$ ($\ell>$ h) has one end connected to A and other to a heavy particle as shown in figure. The particle is held at the level of A with the string tight and released from rest. The height above the plane, where particle is again instantaneously at rest is S, then which is \textbf{correct} ?\par
\qfig{p1_phy_s0_q04_2}
\begin{qoptions}
\opt{A}{velocity of particle on the surface is $\sqrt{\dfrac{2\mathrm{gh}^2}{\ell}}$}
\opt{B}{velocity of particle when it leaves surface $\sqrt{\dfrac{2\mathrm{gh}^4}{\ell^3}}$}
\opt{C}{S = $\dfrac{\mathrm h^5}{\ell^4}$}
\opt{D}{None of these}
\end{qoptions}
\qn{5} ABCDECA is a planar body of mass m of uniform thickness and same material. The dimensions are as shown in the figure. The moment of inertia of the body about an axis passing through point A and perpendicular to planar body is I$_1$ and that of about an axis passing through C and perpendicular to planer body is I$_2$. If I$_1$/I$_2$ is k. Find the value of k.\par
\qfig{p1_phy_s0_q05_1}
\begin{qoptions}
\opt{A}{1}
\opt{B}{2}
\opt{C}{3}
\opt{D}{4}
\end{qoptions}
\qn{6} A block of mass m is projected with velocity v as shown in figure. The ground is smooth but there is friction between A and B. If collision is elastic then choose correct option : (block B is sufficiently long)\par
\qfig{p1_phy_s0_q06_1}
\begin{qoptions}
\opt{A}{the final common velocity of A and B is $\dfrac{2\mathrm v}3$}
\opt{B}{the final common velocity of A and B is $\dfrac{4\mathrm v}9$}
\opt{C}{Energy dissipated in friction is $\dfrac49\mathrm{mv}^2$}
\opt{D}{Energy dissipated in friction is $\dfrac4{27}\mathrm{mv}^2$}
\end{qoptions}
\qn{7} A smooth hemisphere of mass m and radius R is at rest. A smooth solid sphere of mass 2m and radius R moving with velocity V$_0$ between two horizontal smooth surfaces separated by a distance slightly greater than 2R as shown in figure. Solid sphere collides with the hemisphere. If coefficient of restitution is $\dfrac12$, then\par
\qfig{p1_phy_s0_q07_1}
\begin{qoptions}
\opt{A}{The speed of hemisphere after collision is V$_0$}
\opt{B}{The speed of solid sphere after collision is $\dfrac{\mathrm V_0}2$}
\opt{C}{The loss in kinetic energy of the system is $\dfrac14\mathrm{mV}_0^2$}
\opt{D}{The final kinetic energy of hemisphere is $\dfrac14$th the initial kinetic energy of sphere.}
\end{qoptions}
\qn{8} A thin hemispherical shell of uniform density and radius R is placed upside down on a rough inclined plane. The shell just starts toppling without slipping. Given that centre of mass of the hemispherical shell lies at a distance of R/2 from its centre, which of the following is true?
\begin{qoptions}
\opt{A}{Flat surface of the shell should be vertical.}
\opt{B}{Centre of mass of the shell should be vertically above the lowermost point of the shell.}
\opt{C}{$\mu>2$}
\opt{D}{$\mu>\dfrac12$}
\end{qoptions}
\qn{9} A point object of mass m is slipping down on a smooth hemispherical body of mass M and radius R. The point object is tied to a wall with an ideal string as shown. At a certain instant, speed of the hemisphere is v and its acceleration is a. Then speed v$_\mathrm p$ and acceleration a$_\mathrm p$ of a particle has value (Assume all the surfaces in contact are frictionless)\par
\qfig{p1_phy_s0_q09_1}
\begin{qoptions}
\opt{A}{v$_\mathrm P$ = v sin 60$^\circ$}
\opt{B}{v$_\mathrm P$ = v}
\opt{C}{a$_\mathrm P$ = a}
\opt{D}{a$_\mathrm P=\sqrt{\left[\left(\dfrac{\mathrm v^2}{\mathrm R}\right)+\left(\mathrm a\dfrac{\sqrt3}2\right)\right]^2+\left(\dfrac{\mathrm a}2\right)^2}$}
\end{qoptions}
\qn{10} Two blocks A and B of the same mass are connected to a light spring and placed on a smooth horizontal surface. B is given velocity v$_0$ (as shown in the figure) when the spring is in natural length. In the subsequent motion.\par
\qfig{p1_phy_s0_q10_1}
\begin{qoptions}
\opt{A}{the maximum of B will be v$_0$}
\opt{B}{as seen from ground, A can move towards right only}
\opt{C}{the spring will have maximum extension when A and B both stop}
\opt{D}{the spring will be at natural length again when B is at rest}
\end{qoptions}
\qn{11} A bobbin of mass m and moment of inertia I relative to its own axis is being pulled towards right along a horizontal surface by the light string tightly wrapped as shown in figure. There is no slipping on the surface throughout the motion.\par
\qfig{p1_phy_s0_q11_1}
\begin{qoptions}
\opt{A}{If string is pulled by horizontal velocity V, then the angular velocity of bobbin is $\omega=\dfrac{\mathrm V}{3\mathrm r}$}
\opt{B}{If string is pulled by horizontal velocity V, then the angular velocity of bobbin is $\omega=\dfrac{\mathrm V}{2\mathrm r}$}
\opt{C}{If string is pulled by horizontal acceleration a, then the tension in string $\mathrm T=\dfrac{\mathrm a}9\left(4\mathrm m+\dfrac{\mathrm I}{\mathrm r^2}\right)$}
\opt{D}{If string is pulled by horizontal acceleration a, then the tension in string $\mathrm T=\dfrac{\mathrm a}9\left(2\mathrm m+\dfrac{\mathrm I}{\mathrm r^2}\right)$}
\end{qoptions}
\qn{12} Two particles connected with massless rod are rotating with constant angular velocity $\omega$ about axis as shown in figure.\par
\qfig{p1_phy_s0_q12_2}
\begin{qoptions}
\opt{A}{Total angular momentum about O in case (a) and (b) is $2\mathrm{mL}^2\sin^2\theta\,\omega\hat k$.}
\opt{B}{Total angular momentum about O is $2\mathrm{mL}^2\sin^2\theta\,\omega\hat k$ in case (a) only.}
\opt{C}{Total angular momentum about O is $2\mathrm{mL}^2\sin^2\theta\,\omega\hat k$ in case (b) only and net torque about O is zero in this case.}
\opt{D}{In case (a) horizontal torque about O is $\mathrm m\omega^2\mathrm L^2\sin2\theta$}
\end{qoptions}
\qn{13} A disc of mass M radius R moves under gravity while unwinding string under non slip condition. Mark the correct statements :\par
\qfig{p1_phy_s0_q13_1}
\begin{qoptions}
\opt{A}{Acceleration of the disc is $\dfrac{2\mathrm g}3$}
\opt{B}{Rate of ($\vec{\mathrm L}$) change of angular momentum of disc $\dfrac{\mathrm{d}\mathrm L}{\mathrm{dt}}$ about a point along string is in $-\hat k$ direction.}
\opt{C}{There is no point in the space about which angular momentum of disc is conserved}
\opt{D}{Rate of change of angular momentum $\dfrac{\mathrm{dL}}{\mathrm{dt}}$ of disc about center of mass in $+\hat k$ direction.}
\end{qoptions}
\qn{14} Two masses A and B connected with an inextensible string of length $\ell$ lie on a smooth horizontal plane. A is given a velocity v m/s along the ground perpendicular to line AB as shown in figure. The tension in string during their subsequent motion is $\dfrac{\mathrm A}{\mathrm B}\left(\dfrac{\mathrm{mv}^2}{\ell}\right)$, then find $\dfrac{\mathrm A}{\mathrm B}$ ?\par
\qfig{p1_phy_s0_q14_1}
\qn{15} The arm O$'$B of the linkage has a clockwise angular velocity of 10 rad s$^{-1}$ in the position shown in the figure. What is the velocity of A in ms$^{-1}$ (Assume A \& B are rotable joints).\par
\qfig{p1_phy_s0_q15_1}
\qn{16} A homogeneous sphere of radius R = 10 cm rotating about its vertical diameter is gently released on a hole of radius r$_1$ = 8.0 cm made in a fixed horizontal slab. Centre of the hole is vertically below the centre of the sphere. The sphere stops in a time interval $\Delta$t$_1$ = 9.0 s. Calculate time interval required to stop the same sphere, if it is placed rotating with the same angular velocity on a hole of radius r$_2$ = 6.0 cm made in another identical horizontal slab.\par
\qfig{p1_phy_s0_q16_1}
\qn{17} A smooth fixed ring is kept on a smooth horizontal surface. From a point P of the ring a particle is projected at an angle $\alpha$ to the radius vector at P. If e is the coefficient of restitution between the ring and the particle and the particle will return to the point of projection after two reflections. If $\cot^2\alpha=\dfrac{\mathrm A}{\mathrm e^{\mathrm B}}+\dfrac{\mathrm C}{\mathrm e^{\mathrm D}}+\dfrac{\mathrm E}{\mathrm e^{\mathrm F}}$ then find $\dfrac{\mathrm A}{\mathrm B}+\dfrac{\mathrm C}{\mathrm D}+\dfrac{\mathrm E}{\mathrm F}$ ?
\qn{18} The shown structure (a rotable disc over perpendicular rod) rolls purely over a horizontal surface. Obtain velocity of position B. (Use : r = 8m, a = 4m \& $\omega$ = 5 rad/s.)\par
\qfig{p1_phy_s0_q18_1}
\raggedcolumns
\papersec{1}
\partsec{2}{CHEMISTRY}
\qn{19} For a typical carnot cyclic process as shown in the graph, select the correct P--T graph for the entire cyclic process.\par
\qfig{p1_chem_s0_q19_1}
\begin{qoptions}
\opt{A}{\qfigi{p1_chem_s0_q19_2}}
\opt{B}{\qfigi{p1_chem_s0_q19_3}}
\opt{C}{\qfigi{p1_chem_s0_q19_4}}
\opt{D}{\qfigi{p1_chem_s0_q19_5}}
\end{qoptions}
\qn{20} Standard entropy of X$_2$, Y$_2$ and XY$_3$ are 60, 40 and 50 JK$^{-1}$ mol$^{-1}$, respectively. Standard enthalpy of atomisation of X$_2$(g) and Y$_2$(g) are 120 KJ/mole \& 180 KJ/mole respectively. Bond enthalpy of X--Y bond (assume XY$_3$ has three X--Y single bonds) is 120 KJ/mole.\par
$\dfrac12$X$_2$(g) + $\dfrac32$Y$_2$(g) $\to$ XY$_3$(g), then true about the reaction is :
\begin{qoptions}
\opt{A}{spontaneous at all temperature below 1250 K}
\opt{B}{spontaneous at all temperature above 800 K}
\opt{C}{spontaneous at all temperature below 750 K}
\opt{D}{spontaneous at all temperature below 1000 K}
\end{qoptions}
\qn{21} Calculate $|\Delta_\mathrm rG^\circ|$ for (NH$_4$Cl, s) at 350 K.\par
Given : $\Delta_\mathrm fH^\circ$(NH$_4$Cl, s) = $-$314.5 kJ/mol\par
$\mathrm S^\circ_{\mathrm N_2(g)}=192\,\mathrm{JK}^{-1}\mathrm{mol}^{-1}$; \quad $\mathrm S^\circ_{\mathrm H_2(g)}=130.5\,\mathrm{JK}^{-1}\mathrm{mol}^{-1}$\par
$\mathrm S^\circ_{\mathrm{Cl}_2(g)}=223\,\mathrm{JK}^{-1}\mathrm{mol}^{-1}$; \quad $\mathrm S^\circ_{\mathrm{NH}_4\mathrm{Cl}(s)}=94.5\,\mathrm{JK}^{-1}\mathrm{mol}^{-1}$; All given data at 300 K.\par
$\Delta_\mathrm rC_\mathrm P$ = $-$20 J/mol-K ; \quad $\ln\left(\dfrac{350}{300}\right)=0.15$
\begin{qoptions}
\opt{A}{377 kJ/mol}
\opt{B}{374 kJ/mol}
\opt{C}{183 kJ/mol}
\opt{D}{315.5 kJ/mol}
\end{qoptions}
\qn{22} Containers A, B, C and D of equal volume V, each having 120g of X$_4$(linear), X$_2$, X, X$_3$(bent) \& T initial as T, 2T, 3T and 4T respectively. Then find container with maximum $\left\{\dfrac{\Delta\mathrm T}{\mathrm P_\mathrm{final}}\right\}$ value when all gases compress adiabatically to $\dfrac{\mathrm V}3$ against constt ext. pressure. (At mass of X = 60 g/mol.) (ignore vibrational degree of freedom).
\begin{qoptions}
\opt{A}{\fbox{A}}
\opt{B}{\fbox{B}}
\opt{C}{\fbox{C}}
\opt{D}{\fbox{D}}
\end{qoptions}
\qn{23} Which of the following contain correct set of parameters which are both amount and path independent parameters.\par
(I) Internal energy \quad (II) Specific enthalpy\par
(III) Heat exchange at constant pressure condition. \quad (IV) Molar heat exchange at constant volume condition\par
(V) Temperature \quad (VI) Work involved.
\begin{qoptions}
\opt{A}{Only II, IV \& V}
\opt{B}{All option except VI}
\opt{C}{All options except III, IV \& VI}
\opt{D}{Only II \& V}
\end{qoptions}
\qn{24} 1 mole of a real gas changes its state from -- A (2 bar, 3L, 100 K) to state -- B (2 bar, 5 L, 200 K) at constant pressure and finally to state -- C (3 bar, 10 L, 300 K). If $\Delta$U$_\mathrm{BC}$ = 110 J and C$_\mathrm{p,m}$ of gas = 3 R = 3 $\times$ 8.3 JK$^{-1}$ mol$^{-1}$ then choose the correct option(s) :
\begin{qoptions}
\opt{A}{W$_\mathrm{AB}$ = 830 J}
\opt{B}{$\Delta$H$_\mathrm{AC}$ = 4600 J}
\opt{C}{$\Delta$U$_\mathrm{AC}$ = 2200 J}
\opt{D}{$\Delta$U$_\mathrm{AC}$ = 1770 J}
\end{qoptions}
\qn{25} If water vapor is assumed to be a perfect gas, molar enthalpy change for vaporisation of 1 mole of water at 1 bar and 100$^\circ$C is 41 kJ/mol. Choose the \textbf{correct} statement (s)\par
(Take R = 8.3 J/mol/K)
\begin{qoptions}
\opt{A}{$\Delta$U$_\mathrm{vaporisation}$ of 1 mole of water at 1 bar and 100$^\circ$C is nearly 37.904 kJ/mol}
\opt{B}{($\Delta$U $\simeq$ $\Delta$H) for conversion of 1 mole of water into ice at 0$^\circ$C}
\opt{C}{$\Delta$H = 0, for the conversion of 1 mol of water into ice at 0$^\circ$C.}
\opt{D}{($\Delta$H $\simeq$ $\Delta$U) for conversion of 1 mole of water into steam 100$^\circ$C}
\end{qoptions}
\qn{26} Which of the following statement(s) is/are incorrect?
\begin{qoptions}
\opt{A}{The quantities E, H and G have the same dimension}
\opt{B}{Gibb's free energy of 10 gm ice at 0$^\circ$C and 1.0 atm is less than the Gibb's free energy of 10 gm water at 0$^\circ$C and 1 atm.}
\opt{C}{$\Delta$S$_\mathrm{sys}$ = 0 for every adiabatic process in a closed system.}
\opt{D}{For every reversible process, in a closed system, $\Delta$S$_\mathrm{sys}=\dfrac{\Delta\mathrm H_\mathrm{sys}}{\mathrm T}$}
\end{qoptions}
\qn{27} Select the correct statement(s).
\begin{qoptions}
\opt{A}{In a reversible process $\Delta$G is always zero in a closed system.}
\opt{B}{In a reversible process $\Delta$S$_\mathrm{univ}$ is always zero in a closed system.}
\opt{C}{In a reversible process $\Delta$S$_\mathrm{sys}$ is always zero in a closed system.}
\opt{D}{In a reversible process $\Delta$S$_\mathrm{sys}$ is always zero in an isolated system.}
\end{qoptions}
\qn{28} Select the correct option(s):
\begin{qoptions}
\opt{A}{q = nC$_\mathrm v$ dT is applicable to all substances during heating/cooling at constant `v'.}
\opt{B}{q = nC$_\mathrm v$ dT is applicable to ideal gas during heating/cooling at constant `v' only.}
\opt{C}{dU = nC$_\mathrm v$ dT is applicable for real gas at constant `v'}
\opt{D}{dU = nC$_\mathrm v$ dT is applicable for ideal gas at constant `v' only.}
\end{qoptions}
\qn{29} Pick out the incorrect statement among the following :
\begin{qoptions}
\opt{A}{Heat absorb during a change of state (P$_1$V$_1$,T$_1$) to (P$_2$V$_2$,T$_2$) in a closed system involving ideal gas must be equal to $\Delta$H.}
\opt{B}{Heat exchanged during a reversible isothermal process between two fixed states is a difinite quantity.}
\opt{C}{Complete conversion of heat into work is not possible in a reversible isothermal expansion of ideal gas.}
\opt{D}{Complete conversion of heat into work is not possible in an irreversible isothermal expansion of ideal gas.}
\end{qoptions}
\qn{30} A cyclic process for 1 mole of monoatomic gas is shown in the P--T diagram. Choose the correct option :\par
\qfig{p1_chem_s0_q30_1}
\begin{qoptions}
\opt{A}{Work done along path A $\to$ B = $-$ R $\times$ 273}
\opt{B}{Change in internal energy along path BC = $\Delta$U = $\dfrac32$R $\times$ 273}
\opt{C}{For the complete cyclic process $\Delta$U = 0, $\Delta$H = 0, W = (R $\times$ 273) + $\left(2.303\ \mathrm R\times273\log\dfrac12\right)$}
\opt{D}{Heat addition along path CA = q$_{\mathrm C\to\mathrm A}$ = $-\dfrac32$R $\times$ 273}
\end{qoptions}
\qn{31} An ideal gas is expanded isothermally from (P$_1$V$_1$) to (P$_2$,V$_2$) against a constant external pressure. Which of the following is/are correct regarding above process -
\begin{qoptions}
\opt{A}{If P$_\mathrm{ext}$ = P$_2$ piston will not stop automatically at the final state.}
\opt{B}{If P$_\mathrm{ext}$ = P$_2$ piston will stop automatically at the final state}
\opt{C}{Magnitude of work during expansion can vary in the range of 0 $\le$ $|$ W$|$ $\le$ P$_2$ (V$_2$--V$_1$)}
\opt{D}{Work will be zero when P$_\mathrm{ext}$ is zero and will be maximum when P$_\mathrm{ext}$ is equal to pressure of system at final state.}
\end{qoptions}
\qn{32} Calculate the magnitude of net work done in the following cycle for 1 mole of an ideal gas in calorie. Where in process BC, PT = const. (R = 2 cal mol$^{-1}$K$^{-1}$)\par
\qfig{p1_chem_s0_q32_1}
\qn{33} Calculate difference in work done in process--I to process--II.\par
[Given : $\pi$ = 3]\par
\qfig{p1_chem_s0_q33_1}
\qn{34} For AB$_2$(s), lattice dissociation enthalpy is 550 kJ mol$^{-1}$ \& $\Delta$H$_\mathrm{Solution}$ = 50 kJ mol$^{-1}$. If hydration enthalpies (in kJ/mole) of B$^-$ \& A$^{2+}$ ions are in 1 : 4 ratio, then $\dfrac{|\Delta\mathrm H_\mathrm{hydration}|}{10}$ value for cation is ?
\qn{35} When 100 mL of 1.0 M HCl was mixed with 100 mL of 1.0 M NaOH in an insulated beaker at constant pressure, a temperature increase of 5.7$^\circ$C was measured for th beaker and its contents (Expt. 1). Because the enthalpy of neutralization of a strong acid with a strong base is a constant ($-$57.0 kJ mol$^{-1}$), this experiment could be used to measure the calorimeter constant. In a second experiment (Expt.2), 100 mL of 1.0 M acetic acid was mixed with 100 mL of 1.0 M NaOH (under identical conditions to Expt.1) where a temperature rise of 5.524 $^\circ$C was measured. If enthalpy of dissociation of CH$_3$COOH is 2 KJ/mole, then find the \% of CH$_3$COOH self dissociated in the solution.\par
(Consider heat capacity of all solutions are 4.2 J g$^{-1}$ K$^{-1}$ and density of all solutions as 1.0 g mL$^{-1}$)
\qn{36} \qfig{p1_chem_s0_q36_1}
For 1 mole of ideal gas, consider the cyclic process ABCDEA as shown in the graph through 5 steps (A--B, B--C, \dots\dots, E--A).\par
Select number of correct statements.\par
(i) At least 3 states have same volume.\par
(ii) For every step $\dfrac{\mathrm V_\mathrm{final}}{\mathrm V_\mathrm{initial}}\neq1$\par
(iii) V$_\mathrm{maximum}$ = V$_\mathrm B$\par
(iv) Heat inflow occurs in 3--steps.\par
(v) Heat outflow occurs in 2--steps.\par
(vi) Among all steps, maximum $|\Delta$U$|$ value is observed in A$\to$B and D$\to$E processes.\par
(vii) $|\mathrm Q_\mathrm{BC}|=|\mathrm Q_\mathrm{CD}|$\par
(viii) $\Delta\mathrm U_\mathrm{BC}=\Delta\mathrm U_\mathrm{CD}$
\raggedcolumns
\papersec{1}
\partsec{3}{MATHEMATICS}
\qn{37} If the pair of lines $x^2+2xy+\lambda y^2=0$ and $\lambda x^2+2xy+y^2=0$ have exactly one line in common then pair of straight lines representing their other lines is :
\begin{qoptions}
\opt{A}{$x^2+2xy-3y^2=0$}
\opt{B}{$3x^2+2xy-y^2=0$}
\opt{C}{$3x^2+8xy-3y^2=0$}
\opt{D}{$3x^2+10xy+3y^2=0$}
\end{qoptions}
\qn{38} Let S is a circle, from point P outside the circle, tangents are drawn to this circle and points of contact being Q and R. Let A be any point (other than Q and R) on this circle. If a, b, c are length of perpendiculars from A to PQ, PR and QR respectively then
\begin{qoptions}
\opt{A}{2c = a + b}
\opt{B}{c$^2$ = a$^2$ + b$^2$}
\opt{C}{2b = a + c}
\opt{D}{c$^2$ = ab}
\end{qoptions}
\qn{39} On the side BC, CA and AB of triangle ABC, semicircles S$_1$,S$_2$,S$_3$ are described externally. If $\ell_1$ be the length of common tangent to S$_2$ and S$_3$, $\ell_2$ be the length of common tangent to S$_3$ and S$_1$, $\ell_3$ be the length of common tangent to S$_1$ and S$_2$. then ($\ell_1\ell_2\ell_3$) is equal to\par
(where s is semiperimeter and $\Delta$ is area of triangle)
\begin{qoptions}
\opt{A}{2$\Delta$s}
\opt{B}{$\Delta$s}
\opt{C}{$\dfrac{2\Delta^2}{\mathrm s}$}
\opt{D}{$\dfrac{\Delta^2}{\mathrm s}$}
\end{qoptions}
\qn{40} Circles are drawn to cut both the circles $x^2+y^2+6x+5=0$ and $x^2+y^2-6y+5=0$ orthogonally. If all such circle will pass through two fixed points ($x_1$, $y_1$) and ($x_2$, $y_2$) then value of ($x_1+x_2+y_1+y_2$) is equal to
\begin{qoptions}
\opt{A}{3}
\opt{B}{$-$3}
\opt{C}{0}
\opt{D}{6}
\end{qoptions}
\qn{41} Through the point P(2, 3) a pair of perpendicular lines are drawn which meet x-axis point A and B. The locus of incentre of $\Delta$PAB is :
\begin{qoptions}
\opt{A}{$x^2+y^2-4x-6y+13=0$}
\opt{B}{$x^2+y^2-4x-6y+9=0$}
\opt{C}{$x^2-y^2-4x-6y+9=0$}
\opt{D}{$x^2-y^2-4x-6y+13=0$}
\end{qoptions}
\qn{42} Consider circle $x^2+y^2+6x=0$ and $x^2+y^2-2x=0$ which of following are true
\begin{qoptions}
\opt{A}{Equation of direct common tangents are $x-\sqrt3y-3=0$ and $x+\sqrt3y-3=0$}
\opt{B}{Triangle formed by common tangents of two circle is an equilateral triangle}
\opt{C}{Length of direct common tangent of two circle is $2\sqrt3$}
\opt{D}{Two circles intersect each other at 2 real and distinct points}
\end{qoptions}
\qn{43} Tangents are drawn to the circle $x^2+y^2=4$ from the point P($-$4, 0). Points of contact of these tangent are T$_1$ and T$_2$. A rhombus PT$_1$ P$'$T$_2$ is completed then correct statements is/are
\begin{qoptions}
\opt{A}{Circumcentre of the triangle PT$_1$T$_2$ is ($-$2, 0)}
\opt{B}{Ratio of the area of $\Delta$PT$_1$P$'$ to that of $\Delta$P$'$T$_1$T$_2$ is 1 : 2}
\opt{C}{Point P$'$ lies on circle $x^2+y^2=4$}
\opt{D}{Area of quadrilateral PT$_1$ P$'$T$_2$ is $6\sqrt3$}
\end{qoptions}
\qn{44} Circles are drawn through P(2, 9) touching the co-ordinate axes.
\begin{qoptions}
\opt{A}{Radius of one of the circles is 17}
\opt{B}{Radius of one of the circles is 5}
\opt{C}{common chord of two such circles is $x+y-11=0$}
\opt{D}{common chord of two such circles is $x+y+11=0$}
\end{qoptions}
\qn{45} Triangle ABC is formed by the lines $7x-y+3=0$, $x+y-3=0$ and $x-3y-31=0$. Three new lines are drawn parallel to BC, CA, AB passing through A, B and C intersecting each other at P, Q and R respectively then which of the following is correct :
\begin{qoptions}
\opt{A}{circumcentre of triangle PQR lies on $3x+y=3$}
\opt{B}{circumcentre of triangle PQR lies on $x-y=9$}
\opt{C}{circumcentre of $\Delta$PQR, orthocentre and incentre of $\Delta$ABC are collinear}
\opt{D}{Area of $\Delta$PQR is 300}
\end{qoptions}
\qn{46} The line L is drawn through the point (1, 1) to intersect line L$_1$ : $x+3y=1$ and L$_2$ : $x+3y=4$ at P and Q respectively. Line perpendicular to the L from P intersects L$_2$ at R and line perpendicular to L from Q intersects L$_1$ at S. If area of parallelogram PSQR is least then which of the following is/are correct :
\begin{qoptions}
\opt{A}{The equation of line L can be $2x+y-3=0$}
\opt{B}{The equation of line L can be $x-2y+1=0$}
\opt{C}{The equation of line L can be $x+2y-3=0$}
\opt{D}{The least area of PSQR is $\dfrac95$}
\end{qoptions}
\qn{47} Let ABCD be a rectangle inscribed in a circle whose one diameter is along the line $y=x$. If co-ordinates of A and B are (0, 2) and (4, 3) respectively then which of the following is/are correct :
\begin{qoptions}
\opt{A}{point C is $\left(\dfrac{21}5,\dfrac{11}5\right)$}
\opt{B}{point C is (6, 4)}
\opt{C}{point D is $\left(\dfrac15,\dfrac65\right)$}
\opt{D}{point D is (4, 6)}
\end{qoptions}
\qn{48} A circle of radius 3 touches the pair of line $\sqrt3y^2-4xy+\sqrt3x^2=0$ in first quadrant at A and B. If C is centre of circle and O is origin then which of the following is correct.
\begin{qoptions}
\opt{A}{point C is $\left(3+\sqrt3,3+\sqrt3\right)$}
\opt{B}{length of AB is $\dfrac{3(\sqrt3+1)}{\sqrt2}$}
\opt{C}{circumcentre of $\Delta$ABC is $\left(\dfrac{3+\sqrt3}2,\dfrac{3+\sqrt3}2\right)$}
\opt{D}{Area of quadrilateral OACB = $9(2+\sqrt3)$}
\end{qoptions}
\qn{49} Let lines OA and OB are represented by $x^2+y^2+4xy=0$ then possible slopes of AB if $\Delta$OAB is an equilateral triangle is (where O is origin)
\begin{qoptions}
\opt{A}{$\dfrac{\sqrt3}2$}
\opt{B}{$-\dfrac{\sqrt3}2$}
\opt{C}{1}
\opt{D}{$-1$}
\end{qoptions}
\qn{50} A triangle ABC is such that a circle passing through vertex C, centroid G touches side AB at B. If AB = 6, BC = 4 and AC = $\ell$ then [$\ell$] is [where [.] is G.I.F.]
\qn{51} Let 3 lines L$_1$, L$_2$, L$_3$ belonging to family $x-2y+6+\lambda(x-y+2)=0$, where $\lambda$ is a parameter, be internal angle bisectors of $\Delta$ABC. If equation of side AB is $x+3y-4=0$, then find value of $\left[\dfrac{4\Delta}{\sum\left(r\cot\dfrac{\mathrm A}2+a\right)}\right]$ is :\par
(where $\Delta$ is area of $\Delta$ABC and [.] is G.I.F.)
\qn{52} If circle $x^2+y^2-(2\lambda+1)x+(\lambda-7)y+10=0$ passes through points A and B. If angle between tangents to circle at A and B is 90$^\circ$ then the sum of all possible values of $\lambda$ is :
\qn{53} If $16\ell^2+9m^2=24\ell m+6\ell+8m+1$ and $mx+\ell y=1$ is tangent to circle S$_1=0$ and the equation of of circle S$_2$=0 which touches circle S$_1=0$ at (0, 0) and passing through (3, 1) is given by $x^2+y^2-ax+by-c=0$ then value of (3a + b + c) is
\qn{54} If $x^2+y^2+2x+4y+1=0$ and $x^2+y^2-2x-4y+1=0$ are the equation of two excircles of a triangle then area of triangle is
\end{multicols}
\end{document}
